\chapter{矢量球谐函数}
\section{矢量Helmholtz方程}
众所周知，球坐标下标量Helmholtz问题
\begin{align*}
	(\nabla^2 + k^2) \Psi(r, \theta, \phi) = 0.
\end{align*}
的解总可以写成关于基函数\(\psi\)的展开
\begin{align*}
	\Psi = \sum_{l=0}^{\infty} \sum_{m=-l}^{l} a_{lm} \psi_{lm}, \quad
	\psi_{lm}(r, \theta, \phi) = z_l(kr) Y_{lm}(\theta, \phi).
\end{align*}
其中\(Y_{lm}(\theta, \phi)\)为球谐函数 (Spherical Harmonics)，满足微分方程
\begin{align*}
	\left( \frac{1}{\sin\theta} \frac{\partial}{\partial\theta} \left( \sin\theta \frac{\partial}{\partial\theta} \right) + \frac{1}{\sin^2\theta} \frac{\partial^2}{\partial\phi^2} \right) Y_{lm} = -l(l+1) Y_{lm}, \quad
	\frac{\partial Y_{lm}}{\partial\phi} = imY_{lm}.
\end{align*}
\(z_l(x)\)为球贝塞尔 (Spherical Bessel) 函数，满足微分方程
\begin{align*}
	\left( \frac{\dd}{\dd x} \left( x^2 \frac{\dd}{\dd x} \right) + x^2 \right) z_l = l(l+1) z_l.
\end{align*}
并总可以写成\(j_l(x)\) (在\(x=0\)有界) 和\(n_l(x)\) (在\(x=0\)发散) 的线性组合.

以上函数可以用超几何函数表示为
\begin{align*}
	Y_{lm}(\theta, \phi) &= \sqrt{\frac{(2l+1)(l-m)!}{4\pi(l+m)!}} (-1)^m e^{im\phi} (\sin\theta)^m \, _2F_1(-l+m, l+m+1; m+1; \sin^2\frac{\theta}{2}). \\
	j_l(x) &= \frac{x^l}{(2l+1)!!} \, _0F_1(; l+\frac{3}{2}; -\frac{x^2}{4}), \quad
	n_l(x) = -\frac{(2l-1)!!}{x^{l+1}} \, _0F_1(; \frac{1}{2}-l; -\frac{x^2}{4}). \\
	_2F_1(a, b;c; z) &= \sum_{k=0}^{\infty} \frac{(a)_k (b)_k}{(c)_k} \frac{z^k}{k!}, \quad
	_0F_1(; b; z) = \sum_{k=0}^{\infty} \frac{1}{(b)_k} \frac{z^k}{k!}.
\end{align*}
其中
\begin{align*}
	(q)_k = q(q+1)\dots(q+k-1) = \frac{\Gamma(q+k)}{\Gamma(q)}.
\end{align*}

根据Sturm-Liouville理论，\(\psi_{lm}\)可以构成\((r, \theta, \phi)\)的完备标量展开基，在角向，\(Y_{lm}(\theta, \phi)\)具有正交完备性
\begin{align*}
	\int Y_{lm}(\theta, \phi) Y_{l'm'}(\theta, \phi) \dd \Omega = \int_0^{2\pi} \int_0^{\pi} Y_{lm}(\theta, \phi) Y_{l'm'}(\theta, \phi) \sin\theta \, \dd \theta \dd \phi
	= \delta_{ll'} \delta_{mm'}.
\end{align*}
在径向，\(z_l(kr)\)也具有正交完备性，正交性方程形式取决于求解区域和边界条件.

以此为基础，将标量Helmholtz方程的解应用于矢量问题
\begin{align*}
	(\nabla^2 + k^2) \psi(\vb{r}) = 0 \quad\rightarrow\quad (\nabla^2 + k^2) \vb{f}(\vb{r}) = 0.
\end{align*}

当\(\psi_{\lambda}\)是完备基时，可以定义线性无关且各自满足矢量Helmholtz方程的基函数
\begin{align*}
	\vb{L}_{\lambda} = \frac{1}{k} \nabla \psi_{\lambda}, \quad
	\vb{M}_{\lambda} = \frac{1}{k} \nabla \times \left( \vb{a} \psi_{\lambda} \right), \quad
	\vb{N}_{\lambda} = \frac{1}{k} \nabla \times \vb{M}_{\lambda} = \frac{1}{k^2} \nabla \times \left( \nabla \times \left( \vb{a} \psi_{\lambda} \right) \right).
\end{align*}
其中\(\vb{a}\)是常矢量.进而矢量解总可以展开为
\begin{align*}
	\vb{f} = \sum_{\lambda} \left( a_{\lambda} \vb{M}_{\lambda} + b_{\lambda} \vb{N}_{\lambda} + c_{\lambda} \vb{L}_{\lambda} \right).
\end{align*}

易见，\(\vb{L}_{\lambda}\)描述纵向模式，\(\vb{M}_{\lambda}, \vb{N}_{\lambda}\)描述横向模式
\begin{align*}
	\nabla \times \vb{L}_{\lambda} = 0, \quad
	\nabla \cdot \vb{M}_{\lambda} = \nabla \cdot \vb{N}_{\lambda} = 0.
\end{align*}
一般而言，当\(\psi\)描述的不是平面波或球面波时，基函数虽然完备，但不正交、不归一.

\section{球坐标的矢量基函数}
在球坐标中，以上关于基函数的限制条件可适当放宽，所得基函数也具有良好的正交性.不难验证，即使\(\uv{r}\)并非常矢量，所构造的基函数依然满足矢量Helmholtz方程，相应形式为
\begin{align*}
	\psi_{lm}&(r, \theta, \phi) = z_l(kr) Y_{lm}(\theta, \phi). \\
	\vb{L}_{lm} &= \frac{1}{k} \nabla \left( z_l(kr) Y_{lm} \right) = z_l'(kr) Y_{lm} \uv{r} + \frac{1}{k} z_l(kr) \nabla Y_{lm}. \\
	\vb{M}_{lm} &= \frac{1}{k} \nabla \times \left( \uv{r} z_l(kr) Y_{lm} \right) = -\frac{1}{k} z_l(kr) \uv{r} \times \nabla Y_{lm}. \\
	\vb{N}_{lm} &= \frac{1}{k} \nabla \times \vb{M}_{lm} = \frac{l(l+1)}{kr} z_l(kr) Y_{lm} \uv{r} + \frac{1}{k} \frac{\dd(krz_l(kr))}{\dd(kr)} \nabla Y_{lm}.
\end{align*}

它们也具有良好的正交性
\begin{align*}
	\int \vb{L}_{lm} \cdot \vb{L}_{l'm'}^* \dd \vb{r} &= \delta_{ll'} \delta_{mm'} \int \left( \frac{l(l+1)}{k^2} |z_l(kr)|^2 + r^2 |z_l'(kr)|^2 \right) \dd r, \\
	\int \vb{M}_{lm} \cdot \vb{M}_{l'm'}^* \dd \vb{r} &= \delta_{ll'} \delta_{mm'} l(l+1) \int \frac{1}{k^2} |z_l(kr)|^2 \dd r. \\
	\int \vb{N}_{lm} \cdot \vb{N}_{l'm'}^* \dd \vb{r} &= \delta_{ll'} \delta_{mm'} l(l+1) \int \frac{1}{k^2} \left( l(l+1) |z_l(kr)|^2 + \left| \frac{\dd(kr z_l(kr))}{\dd(kr)} \right|^2 \right) \dd r. \\
	\int \vb{L}_{lm} \cdot \vb{M}_{l'm'}^* \dd \vb{r} &= \int \vb{M}_{lm} \cdot \vb{N}_{l'm'}^* \dd \vb{r} = 0, \quad
	\int \vb{L}_{lm} \cdot \vb{N}_{l'm'}^* \dd \vb{r} = \delta_{ll'} \delta_{mm'} \frac{l(l+1)}{k^2}  \left( r |z_l(kr)|^2 \right)_S.
\end{align*}
其中下标\(S\)指在边界上取值.

除去径向部分，角向分布函数可写成如下三种形式
\begin{align*}
	Y_{lm} \uv{r}, \quad
	\vb{X}_{lm} = -\frac{i \vb{r} \times \nabla Y_{lm}}{\sqrt{l(l+1)}}, \quad
	\uv{r} \times \vb{X}_{lm} = \frac{i}{\sqrt{l(l+1)}} r \nabla Y_{lm}.
\end{align*}
即称为矢量球谐函数 (VSH, Vector Spherical Harmonics).

选取球坐标系对应的笛卡尔坐标系，转换关系为
\begin{align*}
	\uv{e}_{\pm} = \frac{\uv{x} \pm i \uv{y}}{\sqrt{2}}, \quad
	\begin{pmatrix} \uv{r} \\ \uv{\theta} \\ \uv{\phi} \end{pmatrix} = \begin{pmatrix}
		\dfrac{\sin\theta}{\sqrt{2}} e^{-i\phi} & \dfrac{\sin\theta}{\sqrt{2}} e^{i\phi} & \cos\theta  \\[8pt]
		\dfrac{\cos\theta}{\sqrt{2}} e^{-i\phi} & \dfrac{\cos\theta}{\sqrt{2}} e^{i\phi} & -\sin\theta \\[8pt]
		-i e^{-i\phi}                           & i e^{i\phi}                            & 0
	\end{pmatrix} \begin{pmatrix} \uv{e}_+ \\ \uv{e}_- \\ \uv{z} \end{pmatrix}.
\end{align*}

进而，在球坐标系中，VSH可以展开为
\begin{align*}
	\vb{X}_{lm} = -\frac{m}{\sqrt{l(l+1)} \sin\theta} Y_{lm} \uv{\theta} - \frac{i}{2} \left( e^{-i\phi} \sqrt{\frac{(l-m)(l+m+1)}{l(l+1)}} Y_{l,m+1} - e^{i\phi} \sqrt{\frac{(l+m)(l-m+1)}{l(l+1)}} Y_{l,m-1} \right) \uv{\phi}.
\end{align*}
其中\(Y_{lm}\uv{r}\)形式普通，\(\uv{r} \times \vb{X}_{lm}\)形式类似，故没有列出.

在直角坐标系中，VSH可以展开为
\begin{align*}
	&\uv{r} Y_{lm} = \left(
	\sqrt{\frac{(l + 1)^2 - m^2}{(2 l + 1) (2 l + 3)}} Y_{l + 1, m} +
	\sqrt{\frac{l^2 - m^2}{4 l^2 - 1}} Y_{l - 1, m}
	\right) \uv{z} \\&\quad+
	\left(
	-\sqrt{\frac{(l + m + 1) (l + m + 2)}{(2 l + 1) (2 l + 3)}} Y_{l + 1, m + 1} +
	\sqrt{\frac{(l - m) (l - m - 1)}{4 l^2 - 1}} Y_{l - 1, m + 1}
	\right) \frac{\uv{e}_-}{\sqrt{2}} \\&\quad-
	\left(
	\sqrt{\frac{(l - m + 1) (l - m + 2)}{(2 l + 1) (2 l + 3)}} Y_{l + 1, m - 1} -
	\sqrt{\frac{(l + m) (l + m - 1)}{4 l^2 - 1}} Y_{l - 1, m - 1}
	\right) \frac{\uv{e}_+}{\sqrt{2}}\\
	&\vb{X}_{lm}=\frac{-i \vb{r}\times \nabla Y_{lm}}{\sqrt{l(l+1)}}=\frac{m}{\sqrt{l(l+1)}}Y_{l,m}\uv{z}+\sqrt{\frac{(l-m+1)(l+m)}{2\,l(l+1)}}Y_{l,m-1}\uv{e}_++\sqrt{\frac{(l+m+1)(l-m)}{2\,l(l+1)}}Y_{l,m+1}\uv{e}_-\\
	&\uv{r}\times\vb{X}_{lm}=i\frac{r\nabla Y_{lm}(\theta,\phi)}{\sqrt{l(l+1)}}=\left(-\sqrt{\frac{l((l+1)^2-m^2)}{(l+1)(2l+1)(2l+3)}}Y_{l+1,m}+\sqrt{\frac{(l+1)(l^2-m^2)}{l(4l^2-1)}}Y_{l-1,m}\right)i\uv{z}\\[6pt]
	&\quad+\left(\sqrt{\frac{l(l+m+1)(l+m+2)}{(l+1)(2l+1)(2l+3)}}Y_{l+1,m+1}+\sqrt{\frac{(l+1)(l-m)(l-m-1)}{l(4l^2-1)}}Y_{l-1,m+1}\right)\frac{i\uv{e}_-}{\sqrt{2}}\\[6pt]
	&\quad-\left(\sqrt{\frac{l(l-m+1)(l-m+2)}{(l+1)(2l+1)(2l+3)}}Y_{l+1,m-1}+\sqrt{\frac{(l+1)(l+m)(l+m-1)}{l(4l^2-1)}}Y_{l-1,m-1}\right)\frac{i\uv{e}_+}{\sqrt{2}}.
\end{align*}

其正交性关系体现在
\begin{align*}
	\left(Y_{l'm'} \uv{r}\right) \cdot \left(\uv{r} \times \vb{X}_{lm}\right) = \left(Y_{l'm'} \uv{r}\right) \cdot \vb{X}_{lm} = 0,\quad
	\int \left(\uv{r} \times \vb{X}_{lm}\right) \cdot \vb{X}_{l'm'}^* \dd \Omega = 0.
\end{align*}

其归一化特性体现在
\begin{align*}
	\int \left(Y_{lm} \uv{r}\right) \cdot \left(Y_{l'm'} \uv{r}\right)^* \dd \Omega =
	\int \left(\uv{r} \times \vb{X}_{lm}\right) \cdot \left(\uv{r} \times \vb{X}_{l'm'}\right)^* \dd \Omega =
	\int \vb{X}_{lm} \cdot \vb{X}_{l'm'}^* \dd \Omega = \delta_{ll'} \delta_{mm'}.
\end{align*}

与之有关的矢量微分公式如：
\begin{align*}
	\nabla \cdot \left(f(r) Y_{lm} \uv{r}\right) = \left(\frac{\dd f}{\dd r} + \frac{2f}{r}\right) Y_{lm},\quad
	\nabla \cdot \left(f(r) \uv{r} \times \vb{X}_{lm}\right) = -\frac{i\sqrt{l(l+1)}}{r} f Y_{lm},
\end{align*}
\begin{align*}
	\nabla \cdot \left(f(r) \vb{X}_{lm}\right) = 0,\quad \nabla \times \left(f(r) \vb{X}_{lm}\right) = \frac{i\sqrt{l(l+1)}}{r} f Y_{lm} \uv{r} + \left(\frac{\dd f}{\dd r} + \frac{f}{r}\right) \uv{r} \times \vb{X}_{lm}.
\end{align*}
\begin{align*}
	\nabla \times \left(f(r) Y_{lm} \uv{r}\right) = -\frac{i\sqrt{l(l+1)} f}{r} \vb{X}_{lm},\quad
	\nabla \times \left(f(r) \uv{r} \times \vb{X}_{lm}\right) = -\left(\frac{\dd f}{\dd r} + \frac{f}{r}\right) \vb{X}_{lm},
\end{align*}

\begingroup
\small
\newpage
\begin{longtable}{ccccc}
	\toprule[1pt]
	\((l,m)\) & 类型                          & \(\uv{e}_+\)                                                                     & \(\uv{e}_-\)                                                                 & \(\uv{z}\)                                                           \\
	\midrule[0.6pt]
	\multirow{3}{*}{\((1,-1)\)}
	          & \(\vb{X}_{lm}\)             & \(0\)                                                                            & \(\sqrt{\dfrac{3}{8\pi}} \cos\theta\)                                        & \(-\sqrt{\dfrac{3}{16\pi}} e^{-i\phi} \sin\theta\)                   \\
	\cmidrule(lr){2-5}
	          & \(\uv{r}\times\vb{X}_{lm}\) & \(-\sqrt{\dfrac{3}{32\pi}} \, i e^{-2i\phi} \sin^2\theta\)                       & \(\sqrt{\dfrac{3}{128\pi}} \, i \left(3+\cos2\theta\right)\)                 & \(-\sqrt{\dfrac{3}{16\pi}} \, i e^{-i\phi} \cos\theta\sin\theta\)    \\
	\cmidrule(lr){2-5}
	          & \(Y_{lm}\uv{r}\)            & \(-\sqrt{\dfrac{3}{16\pi}} e^{-2i\phi} \sin^2\theta\)                            & \(\sqrt{\dfrac{3}{16\pi}} \sin^2\theta\)                                     & \(\sqrt{\dfrac{3}{8\pi}} e^{-i\phi} \cos\theta\sin\theta\)           \\
	\midrule
	\multirow{3}{*}{\((1,0)\)}
	          & \(\vb{X}_{lm}\)             & \(\sqrt{\dfrac{3}{16\pi}} e^{-i\phi} \sin\theta\)                                & \(-\sqrt{\dfrac{3}{16\pi}} e^{i\phi} \sin\theta\)                            & \(0\)                                                                \\
	\cmidrule(lr){2-5}
	          & \(\uv{r}\times\vb{X}_{lm}\) & \(-\sqrt{\dfrac{3}{16\pi}} \, i e^{-i\phi} \cos\theta\sin\theta\)                & \(-\sqrt{\dfrac{3}{16\pi}} \, i e^{i\phi} \cos\theta\sin\theta\)             & \(\sqrt{\dfrac{3}{8\pi}} \, i \sin^2\theta\)                         \\
	\cmidrule(lr){2-5}
	          & \(Y_{lm}\uv{r}\)            & \(-\sqrt{\dfrac{3}{8\pi}} e^{-i\phi} \cos\theta\sin\theta\)                      & \(\sqrt{\dfrac{3}{8\pi}} e^{i\phi} \cos\theta\sin\theta\)                    & \(\sqrt{\dfrac{3}{4\pi}} \cos^2\theta\)                              \\
	\midrule
	\multirow{3}{*}{\((1,1)\)}
	          & \(\vb{X}_{lm}\)             & \(\sqrt{\dfrac{3}{8\pi}} \cos\theta\)                                            & \(0\)                                                                        & \(-\sqrt{\dfrac{3}{16\pi}} e^{i\phi} \sin\theta\)                    \\
	\cmidrule(lr){2-5}
	          & \(\uv{r}\times\vb{X}_{lm}\) & \(-\sqrt{\dfrac{3}{128\pi}} \, i \left(3+\cos2\theta\right)\)                    & \(\sqrt{\dfrac{3}{32\pi}} \, i e^{2i\phi} \sin^2\theta\)                     & \(\sqrt{\dfrac{3}{16\pi}} \, i e^{i\phi} \cos\theta\sin\theta\)      \\
	\cmidrule(lr){2-5}
	          & \(Y_{lm}\uv{r}\)            & \(\sqrt{\dfrac{3}{16\pi}} \sin^2\theta\)                                         & \(-\sqrt{\dfrac{3}{16\pi}} e^{2i\phi} \sin^2\theta\)                         & \(-\sqrt{\dfrac{3}{8\pi}} e^{i\phi} \cos\theta\sin\theta\)           \\
	\midrule
	\multirow{3}{*}{\((2,-2)\)}
	          & \(\vb{X}_{lm}\)             & \(0\)                                                                            & \(\sqrt{\dfrac{5}{8\pi}} e^{-i\phi} \cos\theta \sin\theta\)                  & \(-\sqrt{\dfrac{5}{16\pi}} e^{-2i\phi} \sin^2\theta\)                \\
	\cmidrule(lr){2-5}
	          & \(\uv{r}\times\vb{X}_{lm}\) & \(-\sqrt{\dfrac{5}{32\pi}} \, i e^{-3i\phi} \sin^3\theta\)                       & \(\sqrt{\dfrac{5}{128\pi}}ie^{i\phi} \left(3+\cos2\theta\right)\sin\theta\)  & \(-\sqrt{\dfrac{5}{16\pi}} \, i e^{-2i\phi} \cos\theta\sin^2\theta\) \\
	\cmidrule(lr){2-5}
	          & \(Y_{lm}\uv{r}\)            & \(-\sqrt{\dfrac{15}{128\pi}} e^{-3i\phi} \sin^3\theta\)                          & \(\sqrt{\dfrac{15}{128\pi}} e^{-i\phi} \sin^3\theta\)                        & \(\sqrt{\dfrac{15}{32\pi}} e^{-2i\phi} \cos\theta\sin^2\theta\)      \\
	\midrule
	\multirow{3}{*}{\((2,-1)\)}
	          & \(\vb{X}_{lm}\)             & \(\sqrt{\dfrac{5}{32\pi}} e^{-2i\phi} \sin^2\theta\)                             & \(\sqrt{\dfrac{5}{128\pi}} \left(1+3\cos2\theta\right)\)                     & \(-\sqrt{\dfrac{5}{16\pi}} e^{-i\phi} \cos\theta \sin\theta\)        \\
	\cmidrule(lr){2-5}
	          & \(\uv{r}\times\vb{X}_{lm}\) & \(-\sqrt{\dfrac{5}{8\pi}} \, i e^{-2i\phi} \cos\theta\sin^2\theta\)              & \(\sqrt{\dfrac{5}{8\pi}} \, i \cos^3\theta\)                                 & \(-\sqrt{\dfrac{5}{16\pi}} \, i e^{-i\phi} \cos2\theta\sin\theta\)   \\
	\cmidrule(lr){2-5}
	          & \(Y_{lm}\uv{r}\)            & \(-\sqrt{\dfrac{15}{16\pi}} e^{-2i\phi} \cos\theta\sin^2\theta\)                 & \(\sqrt{\dfrac{15}{16\pi}} \cos\theta\sin^2\theta\)                          & \(\sqrt{\dfrac{15}{8\pi}} e^{-i\phi} \cos^2\theta\sin\theta\)        \\
	\midrule
	\multirow{3}{*}{\((2,0)\)}
	          & \(\vb{X}_{lm}\)             & \(\sqrt{\dfrac{15}{16\pi}} e^{-i\phi} \cos\theta \sin\theta\)                    & \(-\sqrt{\dfrac{15}{16\pi}} e^{i\phi} \cos\theta \sin\theta\)                & \(0\)                                                                \\
	\cmidrule(lr){2-5}
	          & \(\uv{r}\times\vb{X}_{lm}\) & \(-\sqrt{\dfrac{15}{16\pi}} \, i e^{-i\phi} \cos^2\theta\sin\theta\)             & \(-\sqrt{\dfrac{15}{16\pi}} \, i e^{i\phi} \cos^2\theta\sin\theta\)          & \(\sqrt{\dfrac{15}{8\pi}} \, i \cos\theta\sin^2\theta\)              \\
	\cmidrule(lr){2-5}
	          & \(Y_{lm}\uv{r}\)            & \(-\sqrt{\dfrac{5}{128\pi}} e^{-i\phi} \left(1+3\cos2\theta\right)\sin\theta\)   & \(\sqrt{\dfrac{5}{128\pi}} e^{i\phi} \left(1+3\cos2\theta\right)\sin\theta\) & \(\sqrt{\dfrac{5}{256\pi}} \left(5\cos\theta+3\cos3\theta\right)\)   \\
	\midrule
	\multirow{3}{*}{\((2,1)\)}
	          & \(\vb{X}_{lm}\)             & \(\sqrt{\dfrac{5}{128\pi}} \left(1+3\cos2\theta\right)\)                         & \(\sqrt{\dfrac{5}{32\pi}} e^{2i\phi} \sin^2\theta\)                          & \(-\sqrt{\dfrac{5}{16\pi}} e^{i\phi} \cos\theta \sin\theta\)         \\
	\cmidrule(lr){2-5}
	          & \(\uv{r}\times\vb{X}_{lm}\) & \(-\sqrt{\dfrac{5}{8\pi}} \, i \cos^3\theta\)                                    & \(\sqrt{\dfrac{5}{8\pi}} \, i e^{2i\phi} \cos\theta\sin^2\theta\)            & \(\sqrt{\dfrac{5}{16\pi}} \, i e^{i\phi} \cos2\theta\sin\theta\)     \\
	\cmidrule(lr){2-5}
	          & \(Y_{lm}\uv{r}\)            & \(\sqrt{\dfrac{15}{16\pi}} \cos\theta\sin^2\theta\)                              & \(-\sqrt{\dfrac{15}{16\pi}} e^{2i\phi} \cos\theta\sin^2\theta\)              & \(-\sqrt{\dfrac{15}{8\pi}} e^{i\phi} \cos^2\theta\sin\theta\)        \\
	\midrule
	\multirow{3}{*}{\((2,2)\)}
	          & \(\vb{X}_{lm}\)             & \(-\sqrt{\dfrac{5}{8\pi}} e^{i\phi} \cos\theta \sin\theta\)                      & \(0\)                                                                        & \(\sqrt{\dfrac{5}{16\pi}} e^{2i\phi} \sin^2\theta\)                  \\
	\cmidrule(lr){2-5}
	          & \(\uv{r}\times\vb{X}_{lm}\) & \(\sqrt{\dfrac{5}{128\pi}} \, i e^{i\phi} \left(3+\cos2\theta\right)\sin\theta\) & \(-\sqrt{\dfrac{5}{32\pi}} \, i e^{3i\phi} \sin^3\theta\)                    & \(-\sqrt{\dfrac{5}{16\pi}} \, i e^{2i\phi} \cos\theta\sin^2\theta\)  \\
	\cmidrule(lr){2-5}
	          & \(Y_{lm}\uv{r}\)            & \(-\sqrt{\dfrac{15}{128\pi}} e^{i\phi} \sin^3\theta\)                            & \(\sqrt{\dfrac{15}{128\pi}} e^{3i\phi} \sin^3\theta\)                        & \(\sqrt{\dfrac{15}{32\pi}} e^{2i\phi} \cos\theta\sin^2\theta\)       \\
	\bottomrule[1pt]
	\caption{矢量球谐函数直角坐标展开式}
	\label{fig:vsh_cartesian_expansion}
\end{longtable}

\newpage
\begin{longtable}{ccccc}
	\toprule[1pt]
	\((l,m)\) & 展开类型                        & \(\uv{r}\)                                                & \(\uv{\theta}\)                                             & \(\uv{\phi}\)                                                \\
	\midrule[0.6pt]
	\multirow{3}{*}{\((1,-1)\)}
	          & \(\uv{r}Y_{lm}\)            & \(\sqrt{\dfrac{3}{8\pi}}e^{-i\phi}\sin\theta\)            & \(0\)                                                       & \(0\)                                                        \\
	\cmidrule(lr){2-5}
	          & \(\vb{X}_{lm}\)             & \(0\)                                                     & \(\sqrt{\dfrac{3}{16\pi}}e^{-i\phi}\)                       & \(-i\sqrt{\dfrac{3}{16\pi}}e^{-i\phi}\cos\theta\)            \\
	\cmidrule(lr){2-5}
	          & \(\uv{r}\times\vb{X}_{lm}\) & \(0\)                                                     & \(i\sqrt{\dfrac{3}{16\pi}}e^{-i\phi}\cos\theta\)            & \(\sqrt{\dfrac{3}{16\pi}}e^{-i\phi}\)                        \\
	\midrule
	\multirow{3}{*}{\((1,0)\)}
	          & \(\uv{r}Y_{lm}\)            & \(\sqrt{\dfrac{3}{4\pi}}\cos\theta\)                      & \(0\)                                                       & \(0\)                                                        \\
	\cmidrule(lr){2-5}
	          & \(\vb{X}_{lm}\)             & \(0\)                                                     & \(0\)                                                       & \(i\sqrt{\dfrac{3}{8\pi}}\sin\theta\)                        \\
	\cmidrule(lr){2-5}
	          & \(\uv{r}\times\vb{X}_{lm}\) & \(0\)                                                     & \(-i\sqrt{\dfrac{3}{8\pi}}\sin\theta\)                      & \(0\)                                                        \\
	\midrule
	\multirow{3}{*}{\((1,1)\)}
	          & \(\uv{r}Y_{lm}\)            & \(-\sqrt{\dfrac{3}{8\pi}}e^{i\phi}\sin\theta\)            & \(0\)                                                       & \(0\)                                                        \\
	\cmidrule(lr){2-5}
	          & \(\vb{X}_{lm}\)             & \(0\)                                                     & \(\sqrt{\dfrac{3}{16\pi}}e^{i\phi}\)                        & \(i\sqrt{\dfrac{3}{16\pi}}e^{i\phi}\cos\theta\)              \\
	\cmidrule(lr){2-5}
	          & \(\uv{r}\times\vb{X}_{lm}\) & \(0\)                                                     & \(-i\sqrt{\dfrac{3}{16\pi}}e^{i\phi}\cos\theta\)            & \(\sqrt{\dfrac{3}{16\pi}}e^{i\phi}\)                         \\
	\midrule
	\multirow{3}{*}{\((2,-2)\)}
	          & \(\uv{r}Y_{lm}\)            & \(\sqrt{\dfrac{15}{32\pi}}e^{-2i\phi}\sin^2\theta\)       & \(0\)                                                       & \(0\)                                                        \\
	\cmidrule(lr){2-5}
	          & \(\vb{X}_{lm}\)             & \(0\)                                                     & \(\sqrt{\dfrac{5}{16\pi}}e^{-2i\phi}\sin\theta\)            & \(-i\sqrt{\dfrac{5}{16\pi}}e^{-2i\phi}\cos\theta\sin\theta\) \\
	\cmidrule(lr){2-5}
	          & \(\uv{r}\times\vb{X}_{lm}\) & \(0\)                                                     & \(i\sqrt{\dfrac{5}{16\pi}}e^{-2i\phi}\cos\theta\sin\theta\) & \(\sqrt{\dfrac{5}{16\pi}}e^{-2i\phi}\sin\theta\)             \\
	\midrule
	\multirow{3}{*}{\((2,-1)\)}
	          & \(\uv{r}Y_{lm}\)            & \(\sqrt{\dfrac{15}{8\pi}}e^{-i\phi}\cos\theta\sin\theta\) & \(0\)                                                       & \(0\)                                                        \\
	\cmidrule(lr){2-5}
	          & \(\vb{X}_{lm}\)             & \(0\)                                                     & \(\sqrt{\dfrac{5}{16\pi}}e^{-i\phi}\cos\theta\)             & \(-i\sqrt{\dfrac{5}{16\pi}}e^{-i\phi}\cos2\theta\)           \\
	\cmidrule(lr){2-5}
	          & \(\uv{r}\times\vb{X}_{lm}\) & \(0\)                                                     & \(i\sqrt{\dfrac{5}{16\pi}}e^{-i\phi}\cos2\theta\)           & \(\sqrt{\dfrac{5}{16\pi}}e^{-i\phi}\cos\theta\)              \\
	\midrule
	\multirow{3}{*}{\((2,0)\)}
	          & \(\uv{r}Y_{lm}\)            & \(\sqrt{\dfrac{5}{64\pi}}\left(1+3\cos2\theta\right)\)    & \(0\)                                                       & \(0\)                                                        \\
	\cmidrule(lr){2-5}
	          & \(\vb{X}_{lm}\)             & \(0\)                                                     & \(0\)                                                       & \(i\sqrt{\dfrac{15}{8\pi}}\cos\theta\sin\theta\)             \\
	\cmidrule(lr){2-5}
	          & \(\uv{r}\times\vb{X}_{lm}\) & \(0\)                                                     & \(-i\sqrt{\dfrac{15}{8\pi}}\cos\theta\sin\theta\)           & \(0\)                                                        \\
	\midrule
	\multirow{3}{*}{\((2,1)\)}
	          & \(\uv{r}Y_{lm}\)            & \(-\sqrt{\dfrac{15}{8\pi}}e^{i\phi}\cos\theta\sin\theta\) & \(0\)                                                       & \(0\)                                                        \\
	\cmidrule(lr){2-5}
	          & \(\vb{X}_{lm}\)             & \(0\)                                                     & \(\sqrt{\dfrac{5}{16\pi}}e^{i\phi}\cos\theta\)              & \(i\sqrt{\dfrac{5}{16\pi}}e^{i\phi}\cos2\theta\)             \\
	\cmidrule(lr){2-5}
	          & \(\uv{r}\times\vb{X}_{lm}\) & \(0\)                                                     & \(-i\sqrt{\dfrac{5}{16\pi}}e^{i\phi}\cos2\theta\)           & \(\sqrt{\dfrac{5}{16\pi}}e^{i\phi}\cos\theta\)               \\
	\midrule
	\multirow{3}{*}{\((2,2)\)}
	          & \(\uv{r}Y_{lm}\)            & \(\sqrt{\dfrac{15}{32\pi}}e^{2i\phi}\sin^2\theta\)        & \(0\)                                                       & \(0\)                                                        \\
	\cmidrule(lr){2-5}
	          & \(\vb{X}_{lm}\)             & \(0\)                                                     & \(-\sqrt{\dfrac{5}{16\pi}}e^{2i\phi}\sin\theta\)            & \(-i\sqrt{\dfrac{5}{16\pi}}e^{2i\phi}\cos\theta\sin\theta\)  \\
	\cmidrule(lr){2-5}
	          & \(\uv{r}\times\vb{X}_{lm}\) & \(0\)                                                     & \(i\sqrt{\dfrac{5}{16\pi}}e^{2i\phi}\cos\theta\sin\theta\)  & \(-\sqrt{\dfrac{5}{16\pi}}e^{2i\phi}\sin\theta\)             \\
	\bottomrule[1pt]
	\caption{矢量球谐函数球坐标展开式}
	\label{fig:vsh_spherical_expansion}
\end{longtable}
\endgroup